\section{Discussion}\label{sec:discussion}

\subsection{What the label budget buys}
The central trade-off is labels against accuracy. On this stream a fully supervised policy turns 27{,}500 labels into at most 0.007 MCC over SLIDE; equivalently, SLIDE retains 99.2\% of the best supervised MCC with 2{,}000 labels and equals the supervised, error-triggered retraining of its own head. Whether that is the right operating point depends on the price of a verified label. An analyst who must open and adjudicate every message cannot sustain full supervision at gateway volume, and the delayed-label baselines in Table~\ref{tab:main} are an upper reference rather than a deployable alternative. Where labels are cheap, for instance from user reports of a mature organisation, Online-TFIDF-SGD is a strong and inexpensive choice: it was the most accurate method overall and had the lowest false-positive rate on LLM-written legitimate mail.

The ablation and the controls narrow the claim. What SLIDE contributes is \emph{when} it buys labels, and that depends on \emph{where} it looks. At an equal budget, random purchase times cost 0.112 MCC and boundary-aware oracle times gain only 0.004 (Table~\ref{tab:controls}). A KS test on the detector score misses most changes (1.5 alarms for four shifts), and KSWIN on the same score raises 15.4 alarms; projecting $\Phi$ onto random directions exposes shifts that the score hides, because the head maps the new region onto confident but wrong scores. The robust objective matters less than its design suggests. Only the adversarial term has a measurable effect (0.003 MCC); the proximal term and replay show none on this stream, where every update warm-starts from the previous head. A score-triggered policy with the same update comes within 0.007 MCC of SLIDE, but at 3.8 times the label cost and with six times the spread across realisations.

\subsection{Deployment considerations}
SLIDE fits a gateway as an asynchronous side-car. Scoring needs one encoder pass per message while the drift test runs once per batch on cached vectors; an alarm becomes a labelling ticket for the last 500 messages, and the update itself takes seconds on a CPU. The surprisal model contributed nothing measurable here and accounts for about 70\% of the inference cost, so a deployment should drop it unless its own data show otherwise. Two operating rules follow from Section~\ref{sec:results}. The reference must be rebuilt only from post-update traffic, or the test compares the new head's regime with the old one and re-alarms. The purchase window must also be long enough to contain the new campaign; at $n_\ell=250$ the paraphrase segment lost 0.028 MCC.

\subsection{Threats to validity}\label{sec:validity}
\paragraph{Construct validity} The human corpora label spam and phishing together as \emph{malicious}; our detector therefore learns unsolicited bulk mail as well as credential phishing, and the human-written portion of the benchmark is not a pure phishing benchmark. The LLM segments, on the other hand, contain phishing only. The authorship probe (0.9996 macro-F1) shows that LLM-written mail is trivially separable from human mail by style, and the 6.8--27.0\% FPR of the transformer heads on LLM-written legitimate messages indicates that part of what they learn in segment~D is style rather than intent. We report this slice instead of hiding it in the stream average. A stream in which legitimate senders also adopt LLM writing would make this confound more costly.

\paragraph{Internal validity} All hyper-parameters were fixed before the stream runs, and every baseline shares SLIDE's representation and head, so the comparison isolates the policy. The encoder was fine-tuned on warm-up data drawn from the same three human corpora as segment~A, so segment~A is close to in-distribution for $\phi$ and the stream tests generalisation mainly from segment~B onward. The fully supervised baselines are not budget-matched: they receive every label, and only the controls of Table~\ref{tab:controls} compare policies at an equal label count. The baselines are also trained with plain cross-entropy, so part of any gap to them may come from the update rule rather than from the trigger, which the KSWIN-Score+SLIDE-update control measures directly. The baseline detectors use fixed, untuned settings (ADWIN $\delta=0.002$; Page--Hinkley threshold 50 with 200 warm-up observations; KSWIN $\alpha=0.005$ with window 500 and statistic size 150), not values selected on the stream; a tuned ADWIN or KSWIN might narrow or widen the gaps. Ten realisations bound the strength of the realisation-level Wilcoxon tests at an adjusted $p$ of 0.0176, and the batch-level tests treat dependent batches as pairs, so we read them as descriptive.

\paragraph{External validity} The stream is a controlled composition of public corpora with abrupt boundaries, not a recording of one organisation's traffic. Abrupt boundaries favour any window-based test, and each boundary also changes the source corpus, so the shift SLIDE detects mixes campaign change with corpus style. Neither gradual or recurring drift nor label noise was studied. The LLM corpora were generated with OpenAI models only, and the paraphrase campaign uses a single T5 paraphraser; a stronger generator--critic attacker \cite{qi2025} that queries the detector would likely do better than T1. Real drift with an unchanged input marginal is invisible to SLIDE and to every label-free test (Proposition~\ref{prop:sw}).

\paragraph{The padding attack} Prepending 400 characters of legitimate text defeated every transformer detector at every time point (Table~\ref{tab:stress}). The cause is structural: with a 128-token window the lure is pushed out of what $\phi$ and $\psi$ read. Segment-wise encoding with max-pooling over segments would close this specific gap, and so would a cheap whole-message model run in parallel; the TF-IDF results suggest that the second option is effective, and we leave its integration to future work.
